\documentclass[11pt]{article}

\usepackage[utf8]{inputenc}
\usepackage{xspace}
\usepackage{xparse}
\usepackage{usi-technical}

\newcommand{\aufbau}{\textsc{Aufbau}\xspace}
\newcommand{\spg}{\textsc{SPG}\xspace}


\title{\aufbau: Semantic Prefix Grammars and a Unified Constraint Model}
\author{Paul Kronlund-Drouault}
\date{Working Draft}
\setshorttitle{\aufbau: Semantic Prefix Grammars}
\setabstract{%
  Semantic prefix grammars (\spg) are a declarative framework for parsing
  \emph{prefixes} under context-dependent semantic constraints: a partial input
  is accepted exactly when both its syntax and its semantics can still be
  completed to a valid whole. This draft presents the constraint model and a
  single, uniform semantic layer. Syntax is an ordinary grammar with binding
  annotations; semantics is recorded as a constraint graph over the evidence
  exposed by bound children. The semantic layer has one operation, \emph{type
  ascription}, discharged by one mechanism, \emph{unification} of terms over the
  grammar's own signature, regular at the leaves. Types are not a separate
  algebra but terms in the same grammar, so context lookup, equality, inclusion,
  and structural matching are all ascriptions. The soundness target is that the
  parser discards a prefix only when no completion can satisfy its constraints.%
}
\setkeywords{prefix parsing, semantic constraints, unification, constraint graphs, type checking}

\begin{document}

\maketitle

\section{Introduction}
\label{sec:introduction}

\paragraph{Problem.}
Building a formal artifact left to right needs guarantees on the partial result:
a half-written term may already have left the language, broken a typing rule, or
become impossible to complete. Many of these constraints are semantic in the
ordinary sense (names must be in scope, declarations extend later contexts,
branches must agree on types) and remain context-dependent even when the
surface syntax is context-free. The semantic layer must therefore be evaluated
\emph{online}, over partial forests, and must never reject a prefix that still
admits a satisfying completion.

\paragraph{This draft.}
We keep syntax ordinary and make semantics a constraint problem. Productions
carry binding annotations that expose child evidence; semantic rules relate that
evidence by ascription. The accumulated ascriptions form a constraint graph
whose closure under unification yields a three-valued verdict (satisfied,
still-live, or lost), and the parser may prune only on \emph{lost}. Types are
not a separate language: \S\ref{sec:constraints} shows that context lookup,
equality, inclusion, and structural matching are one primitive, ascription by
unification, over terms in the same grammar.

\paragraph{Lineage.}
The mechanism is first-order unification~\cite{robinson1965,martelli1982}, read
as constraint-based type inference~\cite{PottierRemy2005} executed incrementally;
its solver is constraint handling~\cite{Sneyers2009} over an equality core in the
style of equality saturation~\cite{Willsey2021,Tiurin2024}, and its equational and
dependent extensions follow standard unification
theory~\cite{baader2001unification,miller1991}. The grammar-as-signature reading
is the tradition of unification-based grammar~\cite{shieber1986unification};
regular leaves and prefix-aware terminals rest on regular-expression
derivatives~\cite{brzozowski1964derivatives,owens2009regex}; the parser is
Earley-style saturation~\cite{earley1970efficient}; and pending evidence at the
right frontier plays the role of a typed hole~\cite{omar2017hazelnut,zhao2024marking}.

\paragraph{Contributions.}
\begin{enumerate}[nosep]
  \item A constraint model for prefix parsing (\S\ref{sec:foundations}): an
        evidence graph, a three-valued evaluator, and the safe-pruning property.
  \item A unified semantic layer (\S\ref{sec:constraints}): one operation
        (ascription), one mechanism (first-order unification over the grammar's
        signature, regular at the leaves), with types as terms in the same grammar.
  \item Worked instances (\S\ref{sec:examples}) recasting the simply-typed
        lambda calculus and an imperative fragment in the unified model.
\end{enumerate}

% Formal core. The \code macro links definitions to their code counterparts.
% ── Code Reference Infrastructure ───────────────────────────────────────────────
% Autogenerated entries in code-refs-gen.tex (generated by
% scripts/gen-code-refs.sh). This file defines \code and the
% reference-section formatting.
%
% Prerequisites (loaded in main.tex preamble):
%   \usepackage{xparse}
%   \usepackage[colorlinks=true,...]{hyperref}
%
% Usage in paper text:
%   \code{def:constraint-domain}
%   → [def:constraint-domain] with hyperlink.

% ── Inline reference ───────────────────────────────────────────────────────────

\newcommand{\code}[1]{%
  \hyperlink{code:#1}{\textnormal{\texttt{[#1]}}}%
}

% ── Reference section ─────────────────────────────────────────────────────────

\newenvironment{codereferences}{%
  \section*{Code References}
  \small
  \begin{description}[
    leftmargin=1em,
    itemindent=0pt,
    font=\normalfont\ttfamily,
    style=nextline
  ]
}{%
  \end{description}
}

% \coderef{label}{construct}{file}{commit}{line-range}{github-url}
\newcommand{\coderef}[6]{%
  \item[\hypertarget{code:#1}{\texttt{#1}}]%
  \begin{minipage}[t]{\dimexpr\textwidth\relax}%
    \raggedright
    \textbf{#2}\hfill\normalfont\small\texttt{#4}\par
    \smallskip
    \ifx\relax#6\relax
      \texttt{#3}\ \textnormal{L#5}%
    \else
      \href{#6}{\texttt{#3}\ \textnormal{L#5}}%
    \fi
    \par\medskip
  \end{minipage}%
}

\section{Foundations}
\label{sec:foundations}

We fix an alphabet $\Sigma$ and write $L \subseteq \Sigma^\ast$ for a language. A
prefix $s$ is \emph{completable} for $L$ when $s s' \in L$ for some
$s' \in \Sigma^\ast$. Prefix parsing accepts exactly the completable prefixes;
the semantic layer of \S\ref{sec:constraints} refines acceptance so that a prefix
survives only when it is completable to a \emph{semantically} valid whole.

\subsection{Grammar}
\label{sec:grammar}

\begin{definition}[Semantic prefix grammar]
  \label{def:spg}
  \code{def:spg}
  A semantic prefix grammar is a tuple $G = (N, T, P, S, \Theta, \mathcal{T},
  \mathcal{B})$ with nonterminals $N$, terminal recognizers $T$ over segment
  text, productions $P$, start symbol $S \in N$, semantic rules $\Theta$
  (\S\ref{sec:constraints}), a partial map $\mathcal{T} : N \rightharpoonup
  \Theta$ attaching a rule to a nonterminal, and binding identifiers
  $\mathcal{B}$. A production for $n$ is a sequence
  $\alpha_1[b_1]\cdots\alpha_m[b_m]$ where each $\alpha_i \in T \cup N$ and each
  $b_i \in \mathcal{B} \cup \{\varepsilon\}$ names the child's evidence for use
  by a rule. We assume every reachable production has some terminal yield.
\end{definition}

Bindings are syntactic annotations only; they do not affect the context-free
structure. We require $G$ \emph{binding-regular}: each $b \in \mathcal{B}$ occurs
in the productions of exactly one nonterminal and exactly once per production, so
a binding names a unique child position. Terminals are matched through a
segment-level interface returning a status in $\{\mathsf{Prefix},
\mathsf{Exact}, \mathsf{Extensible}\}$ together with a regular residual; this is
the regular-expression derivative~\cite{brzozowski1964derivatives,owens2009regex}
of the recognizer after the consumed segment. The parser depends only on this
status, not on how recognizers are implemented.

\subsection{Evidence}
\label{sec:evidence}

Parsing a prefix yields an AST whose nodes carry status and, at bound positions,
\emph{evidence}: the summary a semantic rule reads.

\begin{definition}[Evidence summary]
  \label{def:evidence}
  \code{def:evidence}
  The summary of an internal node $\langle n, p, (c_1,\dots,c_k)\rangle$ over
  segment span $[i,j)$ is $\nu = \langle n, [i,j), \rho, \tau, \nabla,
  \beta\rangle$, where $\rho \in \{\mathsf{Exact}, \mathsf{Prefix},
  \mathsf{Extensible}\}$ is the status propagated bottom-up (a node is
  $\mathsf{Exact}$ when complete with an exact last child, $\mathsf{Extensible}$
  when complete but the last child may still grow, $\mathsf{Prefix}$ otherwise),
  $\tau$ is its \emph{type term} (\S\ref{sec:constraints}), $\nabla$ is an
  exported context effect (\S\ref{sec:effects}), and $\beta$ resolves bindings.
  For $b$ at child position $i_b$,
  \[
    \beta(\nu, b) =
    \begin{cases}
      \mathsf{Resolved}(\nu_{c_{i_b}}) & i_b \le k, \\
      \mathsf{Pending}(i_b)            & i_b > k,
    \end{cases}
  \]
  so a binding past the right frontier is $\mathsf{Pending}$ until the dot
  reaches it, then $\mathsf{Resolved}$, never the reverse. This is the
  available/unavailable split of incremental attribute
  grammars~\cite{demers1981incremental} and the typed-hole reading of an unfilled
  position~\cite{omar2017hazelnut}.
\end{definition}

\subsection{The constraint graph}
\label{sec:graph}

Semantic information is a graph over evidence, not a tree. Rules relate evidence
summaries directly: the relations need not follow the AST shape.

\begin{definition}[Constraint graph]
  \label{def:constraint-graph}
  \code{def:constraint-graph}
  At a reachable prefix $s$, the constraint graph is a finite hypergraph
  $\mathcal{G}(s) = (\mathcal{E}(s), \mathcal{C}(s))$ whose vertices
  $\mathcal{E}(s)$ are the evidence summaries, terminal leaves, and
  binding-references at $s$, and whose edges $\mathcal{C}(s)$ are the ascription
  constraints emitted so far by fired rules (\S\ref{sec:constraints}).
\end{definition}

A vertex is \emph{frontier} when it still carries an incomplete-input signal:
status in $\{\mathsf{Prefix}, \mathsf{Extensible}\}$, or a $\mathsf{Pending}$
binding-reference. By the propagation rule only the rightmost descendant chain
can be frontier.

\begin{definition}[Open / closed]
  \label{def:open-closed}
  \code{def:open-closed}
  $\mathcal{G}(s)$ is \emph{open} when some edge is incident to a frontier
  vertex, and \emph{closed} when it is not open and the constraint domain finds
  no contradiction in it.
\end{definition}

\subsection{Evaluation and safe pruning}
\label{sec:eval}

\begin{definition}[Constraint domain]
  \label{def:constraint-domain}
  \code{def:constraint-domain}
  A constraint domain is a tuple $D = (\Theta, \mathsf{Closed}, \mathsf{Ctx},
  \mathit{eval}, \oplus)$: a decidable rule language $\Theta$, the closed graphs
  $\mathsf{Closed}$ it admits, semantic contexts $\mathsf{Ctx}$ with empty
  context $\Gamma_0$, the evaluator $\mathit{eval}$ below, and a right-effect
  application $\oplus : \mathsf{Ctx} \times \nabla \to \mathsf{Ctx}$
  (\S\ref{sec:effects}).
\end{definition}

\begin{definition}[Witness, denotation, verdict]
  \label{def:eval}
  \code{def:eval}
  A \emph{witness} for $\mathcal{G}(s)$ is $r \in \Sigma^\ast$ with
  $\mathcal{G}(s r)$ closed; the \emph{denotation} $\mathcal{D}(\mathcal{G}(s))$
  is the set of witnesses. The ideal evaluator is three-valued:
  \[
    \mathit{eval}(\mathcal{G}(s)) =
    \begin{cases}
      \mathsf{Satisfied} & \mathcal{G}(s) \in \mathsf{Closed}, \\
      \mathsf{Live}      & \mathcal{G}(s) \notin \mathsf{Closed}
                           \ \land\ \mathcal{D}(\mathcal{G}(s)) \neq \emptyset, \\
      \mathsf{Lost}      & \mathcal{D}(\mathcal{G}(s)) = \emptyset.
    \end{cases}
  \]
\end{definition}

\begin{lemma}[Safe pruning]
  \label{lem:safe-pruning}
  \code{lem:safe-pruning}
  If the parser discards a prefix $s$ only when $\mathit{eval}(\mathcal{G}(s)) =
  \mathsf{Lost}$, it never discards a prefix with a satisfying completion.
\end{lemma}

\begin{proof}
  $\mathit{eval}(\mathcal{G}(s)) = \mathsf{Lost}$ means
  $\mathcal{D}(\mathcal{G}(s)) = \emptyset$ by Definition~\ref{def:eval}: no
  extension of $s$ closes the graph, so discarding $s$ removes nothing
  completable.
\end{proof}

Realizing this evaluator requires that the implementation's verdict track the
ideal one as input grows. That rests on monotonicity (a hole's completion set
only shrinks under extension, Lemma~\ref{lem:hole-monotone}), so a
$\mathsf{Lost}$ verdict, once reached, is stable. We develop the operational
side in \S\ref{sec:constraints} and state the realizability target in
\S\ref{sec:examples}.

\section{Types, Unification, and Constraints}
\label{sec:constraints}

There is one kind of object, a \emph{term} over the grammar's signature; one
semantic operation, \emph{ascription} (a check of a term against an expected
one); and one mechanism, \emph{unification}. Context lookup, type equality,
inclusion, and structural matching are all ascription discharged by unification.
The type constructors are the grammar's own productions, so there is no separate
algebra of types.

\subsection{Types as terms}

\begin{definition}[Type as term]
  \label{def:term-language}
  \code{def:term-language}
  The productions of the designated type fragment $\mathsf{Ty} \in N$ (written
  \texttt{(*)} on its production in the concrete grammar) form a
  \emph{signature} $\Sigma$: each production is a constructor whose arity is its
  right-hand nonterminals. A \emph{type} is a term in $T(\Sigma)$, a parse tree
  over $\mathsf{Ty}$. There is no privileged type algebra: $\to$, products, and
  unions are constructors only because the grammar writes them. Terminal leaves
  (type names) are recognized by regular terminals, so a type term is a tree
  with \emph{regular leaves}.
\end{definition}

\begin{remark}[The grammar is an order-sorted signature]
  \label{rem:sorts}
  The nonterminals are \emph{sorts}. An alternative whose body is a single
  nonterminal, $n \to n'$, is the subsort declaration $n' \leq n$, and a
  transparent production (Definition~\ref{def:normal-form}) is the injection it
  denotes, which is why collapsing it preserves meaning. The signature is
  therefore order-sorted~\cite{goguen1992order}: a hole at a position of sort
  $n$ ranges over terms of sort at most $n$
  (Definition~\ref{def:hole}), and a term \emph{is a type} exactly when its
  sort is below $\mathsf{Ty}$. This answers, inside the model, which terms may
  stand in type position; no construct outside the fragment can. Order-sorted
  unification stays unitary when sorts with a common subsort have a greatest
  one; this regularity condition is checkable at load time and belongs to the
  realizability gate.
\end{remark}

\subsection{Terms with holes}
\label{sec:patterns}

\begin{definition}[Hole]
  \label{def:hole}
  A \emph{hole} $?A, ?B, \dots$ is a variable standing for an unknown subterm. A
  hole at an internal position ranges over terms; a hole at a leaf position
  ranges over that leaf's regular language, so holes are sorted by the grammar's
  nonterminals (order-sorted, \cite{baader2001unification}). Two are
  distinguished: $\top$, a bare unconstrained hole, the most general type; and
  $\bot$, the absence of any term, a unification failure.
\end{definition}

\begin{definition}[Pattern]
  \label{def:pattern}
  \code{def:pattern}
  A \emph{pattern} is a term with holes, an element of $T(\Sigma, X)$. It is
  \emph{closed} when it has no holes (an ordinary type) and \emph{linear} when no
  hole name repeats. Structure lives in the tree; a repeated hole name is the
  only carrier of equality between positions.
\end{definition}

\subsection{Normal forms}

Equality is compared on normal forms, not on raw trees.

\begin{definition}[Normal form]
  \label{def:normal-form}
  A confluent, terminating rewrite $\mathrm{nf}$ sends each term to a unique
  normal form. The core assumes only the free part: collapse every
  \emph{transparent} node, a production that is delimiters plus a single
  semantic child. Redundant grouping then disappears, so $\mathrm{Int} \to
  (\mathrm{Int} \to \mathrm{Int})$ and $\mathrm{Int} \to \mathrm{Int} \to
  \mathrm{Int}$ are one term. Richer normalizers (declared equations
  \S\ref{sec:ambiguity}, or $\beta$-conversion for a dependent fragment) are
  named extensions; only confluence and termination are assumed here.
\end{definition}

\subsection{Unification}
\label{sec:unify}

\begin{definition}[Unification]
  \label{def:match}
  \code{def:match}
  $\mathrm{unify}(s, t)$ is the most general substitution $\theta$, if one
  exists, with $\theta\,\mathrm{nf}(s) = \theta\,\mathrm{nf}(t)$: structural
  equality of normal forms, leaf holes meeting their regular constraints. This is
  first-order syntactic unification~\cite{robinson1965,martelli1982}: $\top$ (a
  bare hole) unifies with everything, $\bot$ (a constructor clash, or an
  occurs-check failure) with nothing. \emph{Matching} is the one-sided case, a
  pattern against a closed type, and is the basis of ascription.
\end{definition}

\begin{lemma}[Unification is unitary and decidable]
  \label{lem:match-finite}
  \code{lem:match-finite}
  With the occurs-check on (finite terms), $\mathrm{unify}(s, t)$ either fails or
  returns a single most general unifier, unique up to renaming, decidably and in
  near-linear time. No branching arises.
\end{lemma}

\noindent The proof is standard~\cite{martelli1982} and deferred. Unitarity is
the point: the disambiguation that a string-matching model must perform does not
occur, because the parse already fixes which subterm each hole binds. Dropping
the occurs-check admits cyclic terms and gives rational-tree unification for
recursive types (\S\ref{sec:ambiguity}).

\subsection{Ascription: the only operation}

\begin{definition}[Ascription]
  \label{def:ascription}
  The single semantic primitive is the \emph{ascription}
  \[
    b : P,
  \]
  read ``the evidence term $\tau_b$ at $b$ unifies with the expected pattern
  $P$.'' Evaluating it emits $\mathrm{unify}(P, \tau_b)$
  (Definition~\ref{def:evidence}) and contributes the resulting substitution to
  the ambient solution. Read bidirectionally, $P$ is the checked type and
  $\tau_b$ the synthesized one.
\end{definition}

Every typing relation is an ascription, with shared holes the sole carrier of
equality:
\begin{itemize}[nosep]
  \item \emph{Equality} $\tau_1 = \tau_2$ is $\tau_1 : \tau_2$.
  \item \emph{Membership} $x \in \Gamma$ is $x : ?A$ solved against $\Gamma$.
  \item \emph{Inclusion} $\sigma \sqsubseteq \tau$ is subsumption: $\sigma$ is an
        instance of $\tau$ (the one-sided match). $\top$ is above all, $\bot$
        below all.
  \item A structural shape such as $\mathrm{Arrow}(?A, ?B)$ has no built-in
        meaning: $\mathrm{Arrow}$ is a constructor because the grammar wrote it,
        and $?A, ?B$ bind its subterms. Unifying $b : \mathrm{Arrow}(?A, ?B)$
        matches the constructor and binds each subterm by position. No
        decomposition step, and no notion of domain or codomain beyond which
        subterm each hole occupies.
\end{itemize}

\begin{remark}[Subtyping is the subsumption order]
  \label{rem:subtype-inclusion}
  A type at any stage is a term, ground or with holes. Refinement is term
  subsumption,
  \[
    \sigma \sqsubseteq \tau \iff \sigma \text{ is an instance of } \tau,
  \]
  so a hole-bearing (partial, unknown) type is \emph{more general} than its
  instances; $\top$ (a bare hole) is the top, $\bot$ (no unifier) the bottom.
  Unification is the meet: $\mathrm{unify}(\sigma, \tau)$ is their most general
  common instance. Regular sets survive at the leaves, where a leaf hole's value
  is the intersection of its regular constraints, so the prefix layer's regex
  machinery is the leaf case of one structural order. This is at once semantic
  subtyping (a more specific type is an instance) and the prefix order of
  incremental parsing (a partial type is more general, and input only
  instantiates it, Lemma~\ref{lem:hole-monotone}).

  The order is used in one direction. A more general (unknown) actual refutes
  but never confirms: a stable failure against a closed expectation forces
  $\mathsf{Lost}$ and stays failed under instantiation, while a success against
  a still-open actual is provisional. Provisionality costs nothing because
  verdicts are recomputed at every prefix: an edge satisfied today is
  re-checked under the further-instantiated actual tomorrow, and nothing
  irreversible leaves a node before it is exact (its right boundary fixed,
  \S\ref{sec:effects}). So the syntactic layer bounds the semantic one without
  fusing them.
\end{remark}

\subsection{Rules as constraint generators}

\begin{definition}[Rule]
  \label{def:rule}
  \code{def:rule}
  A rule attached to $n \in N$ is a pair $(\Phi, c)$ where the premises
  $\Phi = (b_1 : P_1, \dots, b_m : P_m)$ and the conclusion $c = (\,:\!P_c)$ are
  ascriptions over a shared hole alphabet. Firing the rule at an evidence node
  $\nu$ emits, into the constraint graph
  (Definition~\ref{def:constraint-graph}), one unification edge per premise plus
  the conclusion edge; holes shared between ascriptions become equality edges
  joining the corresponding subterms.
\end{definition}

A rule contributes no control flow of its own: it only \emph{adds edges}, and
shared holes are the joints. This is the constraint-generation reading of type
inference~\cite{PottierRemy2005}: a program is elaborated into a constraint, and
typing is solving it.

\subsection{Solving}

\begin{definition}[Solve]
  \label{def:solve}
  \code{def:solve}
  The solver closes the constraint graph under unification to a fixed point:
  select a pending edge, unify its endpoints, propagate the substitution to
  incident edges, and merge equal holes, until no edge can fire. Closure is
  monotone (substitutions only instantiate, equalities only merge), so the fixed
  point exists and is order-independent. The substrate is union-find over the
  term structure; constraint handling~\cite{Sneyers2009} drives the schedule, and
  the equational extension of \S\ref{sec:ambiguity} represents alternatives as an
  e-graph in the sense of equality saturation~\cite{Willsey2021}.
\end{definition}

\paragraph{Operational closure.}
The implementation realizes the closure bottom-up. A rule's holes are
$\alpha$-renamed at each node, so two firings never collide; the bindings a
node's solution fixes on \emph{foreign} variables (a child's, or one reaching
it through the context) are exported as the \emph{residual equations} of its
evidence, and a node first merges its children's equations into its own
substitution. A binding discovered deep in one subtree thus reaches every
other occurrence of the same metavariable at their least common ancestor:
closure of the constraint graph under unification, computed incrementally.
Metavariable identity is global. A context entry whose type is still a hole
stays one unknown across all its uses; only an explicit instantiation
($\mathsf{inst}$, the elimination of a polymorphic scheme's $\forall$)
freshens. Freshening an ordinary lookup instead would generalize the unknown
silently and accept programs with no type.

The three-valued verdict of \S\ref{sec:guarantees} is read off the closed graph:
\textsf{Satisfied} when every edge unifies and no frontier remains, \textsf{Lost}
when some edge fails \emph{stably}, and \textsf{Live} otherwise. A failure is
stable when no instantiation can repair it. Every failure of the free theory is
stable (a constructor clash or an occurs-violation survives substitution), but
modulo a rewrite theory a hole can block a redex, so normalization and
substitution no longer commute on open terms; there only a failure between
ground terms is stable, and a non-ground failure is postponed as a pending
constraint rather than pruned.

\subsection{Context and effects}
\label{sec:effects}

A rule may extend the context in two ways, and they are kept apart. A premise
\emph{setting} $\Gamma[b:P]$ is \emph{premise-local}: it extends only the
context under which that premise's own term is checked (the descent into the
child the premise binds), and it never reaches a sibling premise. The binder $b$
is itself such a child, resolved under the ambient context, so the setting does
not apply to the descent into $b$. A \emph{right-bound effect} $\nabla$, declared
by the conclusion, is the only thing that crosses to later siblings: it extends
the context they see; a declaration adding a binding is the canonical case. Only
an \emph{exact} node exports an effect, since a prefix node's right boundary is
not yet fixed, and effects compose left to right, applied by $\oplus$
(Definition~\ref{def:constraint-domain}). Operationally this needs two contexts
per node, not one: the context the node was \emph{predicted at} (its
resumption identity, fixed once) and the \emph{working} context that
accumulates left-sibling effects and feeds the next child. Keeping them distinct
is what lets an exported effect reach right siblings without moving where the
node itself closes.

\subsection{Monotonicity and realizability}
\label{sec:guarantees}

\begin{lemma}[Completion monotonicity]
  \label{lem:hole-monotone}
  \code{lem:hole-monotone}
  For an evidence node $e$, the set $\mathrm{comp}(e, s)$ of ground instances of
  its type term only shrinks under extension: $\mathrm{comp}(e, sr) \subseteq
  \mathrm{comp}(e, s)$.
\end{lemma}
\begin{proof}
  Reading input only instantiates holes and grows the context, never the
  reverse: substitution is monotone, so a bound hole stays bound, and a leaf
  hole's regular residual only shrinks under its
  derivative~\cite{owens2009regex}.
\end{proof}

By contraposition $\mathsf{Lost}$ is stable: once no instance satisfies a
constraint, none ever will.

\begin{theorem}[Solver correctness]
  \label{thm:solver-correct}
  \code{thm:solver-correct}
  For the syntactic core the solver verdict on $\mathcal{G}(s)$ equals
  $\mathit{eval}(\mathcal{G}(s))$ (Definition~\ref{def:eval}) at every reachable
  $s$.
\end{theorem}

\noindent\emph{Proof (sketch; full proof deferred).} Induct on the AST. Each
non-$\mathsf{Lost}$ ascription has a satisfying continuation by productivity, and
$\mathsf{Lost}$-stability (Lemma~\ref{lem:hole-monotone}) forbids recovery;
unitarity (Lemma~\ref{lem:match-finite}) makes the per-node unifier unique, so the
join over children matches the open/closed split of
Definition~\ref{def:open-closed}. The equational extension of
\S\ref{sec:ambiguity} is the deferred case.

\subsection{Where multiplicity lives: beyond syntactic unification}
\label{sec:ambiguity}

The syntactic core is unitary, so rule firing never branches and there is no
disambiguation problem: the parse fixes the structure each hole binds.
Multiplicity reappears only when unification runs \emph{modulo} a theory the
grammar declares, each an opt-in extension rather than engine machinery:
\begin{enumerate}[nosep]
  \item \textbf{Equational equivalences.} Declaring $A \mid B \equiv B \mid A$
        and the like makes equality unification modulo an equational
        theory~\cite{baader2001unification}, which is finitary but not unitary
        (associative-commutative unification is the classic
        case~\cite{stickel1981}). The e-graph of Definition~\ref{def:solve}
        represents the alternatives without committing.
  \item \textbf{Recursive types.} Dropping the occurs-check admits cyclic terms;
        unification becomes rational-tree
        unification~\cite{colmerauer1982,Tiurin2024}, still decidable.
  \item \textbf{Conversion.} For a dependent fragment the normal form of
        Definition~\ref{def:normal-form} is $\beta$-normal form and equality is
        conversion, lifting unification to higher-order pattern
        unification~\cite{miller1991}. Decidability rests on strong
        normalization. This is the regime of a small dependent
        kernel~\cite{harper1993lf} and the eventual target of the framework.
\end{enumerate}
The core stays first-order and unitary; the cost is exactly the equational power
a grammar asks for.

\section{Examples}
\label{sec:examples}

\subsection{Simply-typed lambda calculus}

The type fragment is part of the grammar, not a separate algebra
(Definition~\ref{def:term-language}):
\[
  \begin{aligned}
    \text{Identifier} & \to \mathcal{R}(\texttt{[A-Za-z\_][A-Za-z0-9]}^\ast)                                                \\
    \text{Type}       & \to \text{Identifier} \mid \text{Type}\;\text{`$\to$'}\;\text{Type} \mid \text{`('}\,\text{Type}\,\text{`)'} \\
    \text{Variable}   & \to \text{Identifier}[x]                                                            & (\textsc{var}) \\
    \text{Lambda}     & \to \text{`$\lambda$'}\,\text{Identifier}[a]\,\text{`:'}\,\text{Type}[t]\,\text{`.'}\,\text{Expression}[e] & (\textsc{lam}) \\
    \text{Application}& \to \text{Expression}[l]\;\text{Expression}[r]                                      & (\textsc{app}) \\
    \text{Expression} & \to \text{Variable} \mid \text{Lambda} \mid \text{Application}
  \end{aligned}
\]

Each rule is a set of ascriptions (Definition~\ref{def:ascription}); holes are
written $?A$:
\[
  \frac{x : ?A \ \text{in}\ \Gamma}{\,:\,?A}\;(\textsc{var})
  \qquad
  \frac{\Gamma[a{:}t] \vdash e : ?B}{\,:\,t \to ?B}\;(\textsc{lam})
  \qquad
  \frac{l : ?A \to ?B \quad r : ?A}{\,:\,?B}\;(\textsc{app})
\]

The \textsc{app} rule has no idea what a function is. The premise $l : ?A \to ?B$
is one constructor pattern: $\to$ is the Arrow constructor the grammar defines,
$?A, ?B$ its two subterm holes. Unifying it against $l$'s type matches the Arrow
node and binds each hole to a subtree (Definition~\ref{def:match}), first-order
and unitary, with no decomposition step and no notion of domain or codomain
beyond which subterm each hole occupies. The shared $?A$ in $r : ?A$ is the
rule's only equality, forcing $r$'s type to unify with that subterm.
\textsc{lam} extends the context for its body (\S\ref{sec:effects}) and concludes
$t \to ?B$; \textsc{var} ascribes the type found for $x$ in $\Gamma$. The arrow
means only what hole-sharing across the three rules makes it mean.

\paragraph{Prefix (live).}
On $\lambda x{:}\mathrm{Int}$ the lambda node is $\mathsf{Prefix}$: the body
binding $e$ is past the frontier, so $\beta(\nu, e) = \mathsf{Pending}$ and the
\textsc{lam} ascription is incident to a frontier vertex. The graph is open and
the verdict is $\mathsf{Live}$; reading $.\,x$ resolves $e$ and closes it.

\paragraph{Prefix (lost).}
On $\lambda f{:}\mathrm{Int}\to\mathrm{Bool}.\ f\,(\lambda y{:}\mathrm{Int}.\ y\,($
the argument of $f$ has committed to the lambda alternative, so any completion of
that subtree has the Arrow constructor at its head, while \textsc{app} requires it
to unify with the leaf $\mathrm{Int}$. An Arrow-headed term against a leaf is a
constructor clash with no unifier, and by Lemma~\ref{lem:hole-monotone} no
extension recovers it: the verdict is $\mathsf{Lost}$ and the parser may discard
the prefix before the closing parenthesis.

\subsection{Imperative declarations}

A declaration $\texttt{let } x : t = e\,;$ exports a right-bound effect
(\S\ref{sec:effects}) extending the context with $x{:}t$, so a following
statement sees $x$ in scope. Because the effect moves only the working context
and not the declaration node's resumption identity, a block
$\{\texttt{ let } x{:}\mathrm{Int}{=}5;\ \texttt{let } y{:}\mathrm{Int}{=}3; \}$
threads $x$, then $x,y$, through successive statements while each statement still
closes at the context it was predicted in.


\section{Limitations and Open Problems}
\label{sec:limitations}

\paragraph{Beyond the syntactic core.}
The syntactic core is unitary, so match ambiguity is not an open problem
(\S\ref{sec:ambiguity}); multiplicity appears only with the equational and
recursive extensions a grammar may opt into, where unification is finitary, not
unitary. Those extensions, and the dependent-fragment conversion regime, are
stated but not developed here.

\paragraph{Coverage.}
The model targets languages whose type fragment is a grammar signature with
regular leaves. Features that need an equational theory (commutative or
record-like types) or second-order reasoning (row polymorphism, dependency) ride
on the named extensions of \S\ref{sec:ambiguity}, and each new use must be
checked for realizability before admission.

\section{Conclusion}
\label{sec:conclusion}

Semantic prefix grammars keep syntax ordinary and make semantics a constraint
graph solved online. Typing is ascription-by-unification over terms in the
grammar's own signature, with $\top$ a bare variable and $\bot$ a unification
failure. The syntactic core is unitary, so disambiguation is not an open
problem; the remaining work is the equational, recursive, and dependent
extensions and the mechanized proof of the semantic invariants.

\section*{Acknowledgments}

We thank Niels M\"undler at ETH Z\"urich's SRI Lab for feedback on the
implementation and theoretical advice. This work was supported by Unsuspicious
Industries, an independent nonprofit research collective.

\bibliographystyle{plain}
\bibliography{references}

% Code-reference appendix omitted while code `///` labels are reconciled with the
% rewritten definitions; regenerate with scripts/gen-code-refs.sh and restore
% \begin{codereferences}% Auto-generated by scripts/gen-code-refs.sh — do not edit.
% Commit: f04c12d  |  Generated: 2026-05-20 17:06:18 UTC

\coderef{def:constraint-domain}{fn apply_effect}{src/semantics/domain.rs}{f04c12d}{157--160}{https://github.com/Unsuspicious-Industries/aufbau/blob/f04c12d65e09699c87695c300d2b3af36868e524/src/semantics/domain.rs#L157-L160}
\coderef{def:eval-impl}{fn finalize}{src/semantics/domain.rs}{f04c12d}{132--146}{https://github.com/Unsuspicious-Industries/aufbau/blob/f04c12d65e09699c87695c300d2b3af36868e524/src/semantics/domain.rs#L132-L146}
\coderef{def:evidence-graph}{#[derive(Clone, Copy, Debug, PartialEq, Eq)]}{src/engine/parse/arena.rs}{f04c12d}{92--104}{https://github.com/Unsuspicious-Industries/aufbau/blob/f04c12d65e09699c87695c300d2b3af36868e524/src/engine/parse/arena.rs#L92-L104}
\coderef{def:evidence-summary}{type Effect}{src/semantics/domain.rs}{f04c12d}{101--103}{https://github.com/Unsuspicious-Industries/aufbau/blob/f04c12d65e09699c87695c300d2b3af36868e524/src/semantics/domain.rs#L101-L103}
\coderef{def:open-closed}{#[derive(Clone, Copy, Debug, PartialEq, Eq)]}{src/engine/parse/arena.rs}{f04c12d}{92--104}{https://github.com/Unsuspicious-Industries/aufbau/blob/f04c12d65e09699c87695c300d2b3af36868e524/src/engine/parse/arena.rs#L92-L104}
\coderef{lem:evidence-monotone}{trait ConstraintDomain}{src/semantics/domain.rs}{f04c12d}{52--72}{https://github.com/Unsuspicious-Industries/aufbau/blob/f04c12d65e09699c87695c300d2b3af36868e524/src/semantics/domain.rs#L52-L72}
\coderef{lem:evidence-realizable}{fn bind_asc}{src/domains/typing/domain.rs}{f04c12d}{129--140}{https://github.com/Unsuspicious-Industries/aufbau/blob/f04c12d65e09699c87695c300d2b3af36868e524/src/domains/typing/domain.rs#L129-L140}
\coderef{lem:prefix-completeness}{fn check_all_prefixes_parseable}{src/validation/parseable/mod.rs}{f04c12d}{105--111}{https://github.com/Unsuspicious-Industries/aufbau/blob/f04c12d65e09699c87695c300d2b3af36868e524/src/validation/parseable/mod.rs#L105-L111}
\coderef{lem:rule-realizable}{trait ConstraintDomain}{src/semantics/domain.rs}{f04c12d}{52--72}{https://github.com/Unsuspicious-Industries/aufbau/blob/f04c12d65e09699c87695c300d2b3af36868e524/src/semantics/domain.rs#L52-L72}
\coderef{lem:safe-pruning}{fn finalize}{src/semantics/domain.rs}{f04c12d}{132--146}{https://github.com/Unsuspicious-Industries/aufbau/blob/f04c12d65e09699c87695c300d2b3af36868e524/src/semantics/domain.rs#L132-L146}
\coderef{lem:typeof-realizable}{fn bind_asc}{src/domains/typing/domain.rs}{f04c12d}{129--140}{https://github.com/Unsuspicious-Industries/aufbau/blob/f04c12d65e09699c87695c300d2b3af36868e524/src/domains/typing/domain.rs#L129-L140}
\coderef{prop:binding-uniqueness}{#[derive(Debug, Clone, Default)]}{src/engine/binding.rs}{f04c12d}{18--21}{https://github.com/Unsuspicious-Industries/aufbau/blob/f04c12d65e09699c87695c300d2b3af36868e524/src/engine/binding.rs#L18-L21}
\coderef{sec:context-flow}{fn apply_effect}{src/semantics/domain.rs}{f04c12d}{157--160}{https://github.com/Unsuspicious-Industries/aufbau/blob/f04c12d65e09699c87695c300d2b3af36868e524/src/semantics/domain.rs#L157-L160}
\coderef{sec:engine-parsing}{fn finalize}{src/semantics/domain.rs}{f04c12d}{132--146}{https://github.com/Unsuspicious-Industries/aufbau/blob/f04c12d65e09699c87695c300d2b3af36868e524/src/semantics/domain.rs#L132-L146}
\coderef{sec:gram-def}{trait ConstraintLoader}{src/semantics/loader.rs}{f04c12d}{17--22}{https://github.com/Unsuspicious-Industries/aufbau/blob/f04c12d65e09699c87695c300d2b3af36868e524/src/semantics/loader.rs#L17-L22}
\coderef{sec:meta-compilation}{#[must_use]}{src/domains/typing/compiler.rs}{f04c12d}{79--89}{https://github.com/Unsuspicious-Industries/aufbau/blob/f04c12d65e09699c87695c300d2b3af36868e524/src/domains/typing/compiler.rs#L79-L89}
\coderef{sec:semantic-domains}{trait ConstraintDomain}{src/semantics/domain.rs}{f04c12d}{52--72}{https://github.com/Unsuspicious-Industries/aufbau/blob/f04c12d65e09699c87695c300d2b3af36868e524/src/semantics/domain.rs#L52-L72}
\coderef{sec:typing-domain}{#[derive(Clone, Debug)]}{src/domains/typing/domain.rs}{f04c12d}{25--25}{https://github.com/Unsuspicious-Industries/aufbau/blob/f04c12d65e09699c87695c300d2b3af36868e524/src/domains/typing/domain.rs#L25-L25}
\coderef{thm:completability-soundness}{fn check_all_prefixes_parseable}{src/validation/parseable/mod.rs}{f04c12d}{105--111}{https://github.com/Unsuspicious-Industries/aufbau/blob/f04c12d65e09699c87695c300d2b3af36868e524/src/validation/parseable/mod.rs#L105-L111}
\coderef{thm:prefix-monotonicity}{fn check_all_prefixes_parseable}{src/validation/parseable/mod.rs}{f04c12d}{105--111}{https://github.com/Unsuspicious-Industries/aufbau/blob/f04c12d65e09699c87695c300d2b3af36868e524/src/validation/parseable/mod.rs#L105-L111}
\coderef{thm:typing-realizable}{fn finalize}{src/semantics/domain.rs}{f04c12d}{132--146}{https://github.com/Unsuspicious-Industries/aufbau/blob/f04c12d65e09699c87695c300d2b3af36868e524/src/semantics/domain.rs#L132-L146}

\end{codereferences}.

\end{document}
